%!TEX program = xelatex
\documentclass[titlepage,oneside]{ctexart}
\usepackage{ecustpc} % 加载刚才定义的样式

% 定义竞赛信息
\title{2026 年秋季华东理工大学程序设计竞赛试题册}
\author{ECUSTPC 2026 Fall Jury}
\date{X, 2026}

\begin{document}

% 设置页眉内容：竞赛名称，日期
\setcontestheader{2026 年秋季华东理工大学程序设计竞赛 正式赛}{2026 年 X 月 X 日}
\pagestyle{fancy}

\setlength{\headheight}{25.3186pt}

\begin{titlepage}
\begin{center}
    \huge{\textbf{2026 年秋季华东理工大学程序设计竞赛}} \\ [10pt]
    \huge{正式赛 \quad 试题册} \\ [10pt]
    \Large{2026 年 X 月 X 日 X$\sim$X}
\end{center}

\vspace{2em}
{\large\BF{试题列表}}

\begin{table}[htbp]
  \begin{center}
    \large
    \begin{tabular}{ccccc}
      \toprule
      题号 & 试题名称 & 时空限制 & 注释 \\
      \midrule
      \myboxi{\BF{\TT{A}}} & \BF{纽约提喻法} & X/X & \\
      \myboxii{\BF{\TT{B}}} & \BF{AIMC 实验室} & X/X & \\
      \myboxiii{\BF{\TT{C}}} & \BF{一分为二 2} & X/X & \\
      \myboxiv{\BF{\TT{D}}} & \BF{单位距离！？} & X/X & \\
      \myboxv{\BF{\TT{E}}} & \BF{占位标题 E} & X/X & \\
      \myboxvi{\BF{\TT{F}}} & \BF{凸包} & X/X & \\
      \myboxvii{\BF{\TT{G}}} & \BF{星之彼端} & X/X & 通信题 \\
      \myboxviii{\BF{\TT{H}}} & \BF{一般攻击魔法} & X/X & \\
      \myboxix{\BF{\TT{I}}} & \BF{天外来物} & X/X & \\
      \myboxx{\BF{\TT{J}}} & \BF{时间谜题 2} & X/X & \\
      \myboxxi{\BF{\TT{K}}} & \BF{永世沙漏} & X/X & \\
      \myboxxii{\BF{\TT{L}}} & \BF{净化行动 3} & X/X & \\
      \myboxiii{\BF{\TT{M}}} & \BF{星光站} & X/X & \\
      \bottomrule
    \end{tabular}
  \end{center}
\end{table}

试题册共包含 13 道试题和 X 页.

若您的试题册缺页请立刻联系志愿者。

请务必认真阅读位于比赛文件附件中的赛事规则和技术规则。

赛事开始前请勿翻阅试题册，亦请勿触碰鼠标和键盘。

请注意黑白打印的试题册可能会丢失部分需要颜色展现的信息。

\textit{写在题目描述前的斜体文字 是 Feature Text（风味文本），和题目无关，包括这句话。}
\end{titlepage}

\linespread{1.0}
\setlength{\parskip}{0\baselineskip}
\normalsize

\newpage
\setcounter{page}{1}

\section[A. 纽约提喻法]{\myboxi{      \BF{\TT{A}}} 纽约提喻法}

\textit{运行时间限制}   X second

\textit{内存使用限制} X Gigabyte

\textit{读入文件} \qquad 标准输入

\textit{输出文件} \qquad 标准输出


\begin{quote}
	\textit{And they say there is no fate, but there is, it's what you create.}

	\textit{And even though the world goes on for eons and eons...}

	\textit{...you are only here for a fraction of a fraction of a second.}
\end{quote}

\subsection{题目描述}

凯顿对"命运"有了自己的理解，他不再借助戏剧，而想用一串字符来提喻自己的人生。然而，当他尝试构建这个提喻时，为了追求纯粹、忠于自我的原创表达，他再度陷入了自我指涉的无限循环。于是这一次，凯顿将选择权交给了你，由你来描绘他的人生。

凯顿的提喻是一个长度为 $L$ 的字符串，记为 $T$。称 $T$ 是\BF{合法}的，当且仅当它同时满足以下三个条件：

\begin{enumerate}
	\item $T$ 仅由小写英文字母组成。
	\item 对于每个 $i = 1, 2, \ldots, K$，字符串 $T$ 的第 $p_i$ 个字符必须等于给定的字符 $c_i$。这里 $p_i$ 是位置下标，满足 $1 \le p_i \le L$，$c_i$ 为小写字母。
	\item 对于每个 $i = 1, 2, \ldots, N$，给定的字符串 $s_i$ \BF{不能}作为 $T$ 的连续子串出现。
\end{enumerate}

现在给定 $L$、$K$ 组位置约束 $(p_i, c_i)$ 以及 $N$ 个禁止子串 $s_i$，请你计算所有合法字符串 $T$ 的总数，并将结果对 $998\,244\,353$ 取模后输出。

\subsection{输入格式}

第一行输入三个整数 $L$, $K$, $N$ ($1 \le L \le 10^9$, $0 \le K \le 10$, $0 \le N \le 50$)，分别表示字符串长度、位置约束数量和禁止子串数量。

随后 $K$ 行，每行输入一个整数 $p_i$ 和一个小写字母 $c_i$ ($1 \le p_i \le L$)，表示一个位置约束。保证所有 $p_i$ 互不相同。

随后 $N$ 行，每行输入一个非空字符串 $s_i$，表示一个禁止子串。

保证所有 $s_i$ 的长度之和不超过 $100$（即 $\sum_{i=1}^N \vert{}s_i\vert{} \le 100$）。

\subsection{输出格式}

对于每组测试数据，输出一行一个整数，表示合法字符串的总数对 $998\,244\,353$ 取模后的结果。

\subsection{样例 1}

\begin{minipage}{\linewidth}
	\begin{minipage}[t]{0.49\linewidth}
		\begin{lstlisting}[title=\texttt{standard input}]
3 1 1
2 b
ab
\end{lstlisting}
	\end{minipage}
	\begin{minipage}[t]{0.49\linewidth}
		\begin{lstlisting}[title=\texttt{standard output}]
650
\end{lstlisting}
	\end{minipage}
\end{minipage}


\subsection{样例 1 解释}

字符串长度为 $3$，第 $2$ 个字符必须是 \TT{b}，且不能包含子串 \TT{ab}。合法的字符串形如 \TT{?b?}，其中第 $1$ 个位置可以是除 \TT{a} 以外的任意小写字母，共 $25 \times 26 = 650$ 种。


\newpage
\section[B. AIMC实验室]{\myboxii{      \BF{\TT{B}}} AIMC实验室}

\textit{运行时间限制}   X second

\textit{内存使用限制} X Gigabyte

\textit{读入文件} \qquad 标准输入

\textit{输出文件} \qquad 标准输出


\begin{quote}
	\textit{街灯璀璨，可我日渐透明。}
\end{quote}

\subsection{题目描述}

你所在的 AIMC 实验室研究出了一种全新的三维量子矩阵！

给定一个大小为 $N \times M \times H$ 的三维空间网格，坐标范围为 $1 \le x \le N$, $1 \le y \le M$, $1 \le z \le H$。我们用 $u$ 代表坐标为 $(x_u, y_u, z_u)$ 的量子。

对于网格中的每一个量子，你可以将其初始化为两种基态之一：$|0\rangle$ 态或 $|1\rangle$ 态。

对于任意两个量子 $u$ 和 $v$，定义它们之间的\BF{平方空间距离}为：
$$
d^2(u, v) = (x_u - x_v)^2 + (y_u - y_v)^2 + (z_u - z_v)^2
$$

系统存在 $8$ 个大小同为 $N \times M \times H$ 的能量矩阵 $A_1, A_2, \ldots, A_8$。系统总能量计算规则如下：

对于网格中的任意量子 $u$：

\begin{itemize}
	\item 若 $u$ 为 $|0\rangle$ 态，总能量增加 $A_1(u)$ ；
	\item 若 $u$ 为 $|1\rangle$ 态，总能量增加 $A_2(u)$ 。
\end{itemize}

定义 $S_1(u)$ 为网格内所有满足 $d^2(u, v) = 1$ 的量子 $v$ 的集合：

\begin{itemize}
	\item 若 $u$ 为 $|0\rangle$ 态，且 $S_1(u)$ 中所有量子均为 $|1\rangle$ 态，总能量增加 $A_3(u)$ ；
	\item 若 $u$ 为 $|1\rangle$ 态，且 $S_1(u)$ 中所有量子均为 $|0\rangle$ 态，总能量增加 $A_4(u)$ 。
\end{itemize}

定义 $S_2(u)$ 为网格内所有满足 $d^2(u, v) = 2$ 的量子 $v$ 的集合：

\begin{itemize}
	\item 若 $u$ 为 $|0\rangle$ 态，且 $S_2(u)$ 中所有量子均为 $|0\rangle$ 态，总能量增加 $A_5(u)$ ；
	\item 若 $u$ 为 $|1\rangle$ 态，且 $S_2(u)$ 中所有量子均为 $|1\rangle$ 态，总能量增加 $A_6(u)$ 。
\end{itemize}

定义 $S_3(u)$ 为网格内所有满足 $d^2(u, v) = 3$ 的量子 $v$ 的集合：

\begin{itemize}
	\item 若 $u$ 为 $|0\rangle$ 态，且 $S_3(u)$ 中所有量子均为 $|1\rangle$ 态，总能量增加 $A_7(u)$ ；
	\item 若 $u$ 为 $|1\rangle$ 态，且 $S_3(u)$ 中所有量子均为 $|0\rangle$ 态，总能量增加 $A_8(u)$ 。
\end{itemize}

作为AIMC实验室的首席科学家，请你为网格中的每个量子指定一个状态，使得系统总能量最大化，并输出这个最大值。

\subsection{输入格式}

第一行输入三个整数 $N$, $M$, $H$ ($1 \le N, M, H \le 20$)，表示三维网格的大小。

接下来依次给出 8 个能量观测矩阵 $A_1$ 到 $A_8$ 的数据。

对于每个矩阵，包含 $N$ 个区块。每个区块包含 $M$ 行，每行 $H$ 个整数。区块之间用空行隔开。

第 $x$ 个区块的第 $y$ 行第 $z$ 个数即代表对应坐标 $(x, y, z)$ 的能量值，所有矩阵中的能量值 $a_i \in [0, 10^9]$。

\subsection{输出格式}

输出一行一个整数，表示系统可以获得的最大总能量值。

\subsection{样例 1}

\begin{minipage}{\linewidth}
	\begin{minipage}[t]{0.49\linewidth}
		\begin{lstlisting}[title=\texttt{standard input}]
1 1 1
3

5

100

200

1000

2000

50

60
\end{lstlisting}
	\end{minipage}
	\begin{minipage}[t]{0.49\linewidth}
		\begin{lstlisting}[title=\texttt{standard output}]
2265
\end{lstlisting}
	\end{minipage}
\end{minipage}


\subsection{样例 1 解释}

整个网格仅有一个量子 $(1, 1, 1)$。

将该量子初始化为 $|1\rangle$ 时获得 $A_2 + A_4 + A_6 + A_8 = 5 + 200 + 2000 + 60 = 2265$，为最优。

\subsection{样例 2}

\begin{minipage}{\linewidth}
	\begin{minipage}[t]{0.49\linewidth}
		\begin{lstlisting}[title=\texttt{standard input}]
2 2 2
1 1
1 1

1 1
1 1

0 0
0 0

0 0
0 0

0 0
0 0

0 0
0 0

0 0
0 0

0 0
0 0

10 10
10 10

10 10
10 10

0 0
0 0

0 0
0 0

0 0
0 0

0 0
0 0

0 0
0 0

0 0
0 0
\end{lstlisting}
	\end{minipage}
	\begin{minipage}[t]{0.49\linewidth}
		\begin{lstlisting}[title=\texttt{standard output}]
88
\end{lstlisting}
	\end{minipage}
\end{minipage}


\subsection{样例 2 解释}

将 $2 \times 2 \times 2$ 中所有 $8$ 个量子统一初始化为 $|0\rangle$。

\begin{itemize}
	\item 每个量子从 $A_1$ 各获得 $1$ 的能量，合计 $8$。
	\item 每个量子距离为 $2$ 的邻居也全部为 $|0\rangle$，从 $A_5$ 各获得 $10$，合计 $80$。
\end{itemize}

总能量 $8 + 80 = 88$。可以证明这是最大值。

\newpage
\section[C. 一分为二 2]{\myboxiii{      \BF{\TT{C}}} 一分为二 2}

\textit{运行时间限制}   2 seconds

\textit{内存使用限制} 1 Gigabyte

\textit{读入文件} \qquad 标准输入

\textit{输出文件} \qquad 标准输出


\begin{quote}
	\textit{别在昨日里频频回首，说什么将寂寞对半分担。}
\end{quote}

\subsection{题目描述}

给定两个长度为 $n$ 的数组 $a$ 和 $b$。

你可以任选 $k \in [0,n]$，将 $a$ 的前 $k$ 个元素和 $b$ 的后 $n-k$ 个元素拼接成一个长度为 $n$ 的新数组。

请从新数组中选出恰好 $m$ 个互不相邻的元素。记其中最小的元素为 $p$，求 $p$ 的最大值。

\subsection{输入格式}

第一行输入一个整数 $T$ ($1 \le T \le 10^4$)，表示测试数据的数量。

每组测试数据第一行输入两个整数 $n$ 和 $m$ ($1 \le m \le \lceil \frac{n}{2} \rceil \le n \le 2 \times 10^5$)，分别表示数组长度和需要选择的元素数量。

第二行输入 $n$ 个整数 $a_1, a_2, \ldots, a_n$ ($-10^9 \le a_i \le 10^9$)，表示数组 $a$。

第三行输入 $n$ 个整数 $b_1, b_2, \ldots, b_n$ ($-10^9 \le b_i \le 10^9$)，表示数组 $b$。

保证所有测试数据的 $\sum n \le 2 \times 10^5$。

\subsection{输出格式}

对于每组测试数据，输出一行一个整数，表示 $p$ 的最大值。

\subsection{样例 1}

\begin{minipage}{\linewidth}
	\begin{minipage}[t]{0.49\linewidth}
		\begin{lstlisting}[title=\texttt{standard input}]	
2
5 2
5 4 3 2 1
1 2 3 4 5
9 5
-1 1 -4 -5 -1 4 -1 -9 1
-9 -9 -8 2 4 -4 -3 -5 3
\end{lstlisting}
	\end{minipage}
	\begin{minipage}[t]{0.49\linewidth}
		\begin{lstlisting}[title=\texttt{standard output}]
5
-4
\end{lstlisting}
	\end{minipage}
\end{minipage}

\subsection{样例 1 解释}

对于第 $1$ 组测试数据，当 $k = 4$ 时，新数组为 $[5, 4, 3, 2, 5]$，选择第 $1,5$ 个元素，最小元素为 $5$，可以证明这是最大值。

\newpage
\section[D. 位距！？]{\myboxiv{      \BF{\TT{D}}} 位距！？}

\textit{运行时间限制}   2 seconds

\textit{内存使用限制} 1 Gigabyte

\textit{读入文件} \qquad 标准输入

\textit{输出文件} \qquad 标准输出


\begin{quote}
	\textit{最后跳支舞吧。}

	\textit{在魔都的远郊，}

	\textit{在化工强校的信息楼，}

	\textit{在这个光怪陆离的}

	\textit{革命前夜。}
\end{quote}

\subsection{题目描述}

请在欧氏平面上构造一个 $n$ 个点的点集，使得：

\begin{itemize}
	\item 点的横纵坐标绝对值不超过 $10^5$；
	\item 任意两点之间的距离不小于 $10^{-4}$；
	\item 至少存在 $\left\lfloor \frac{1}{2}n\log_2 n - n \right\rfloor$ 个不同的无序点对距离在 $[1 - 10^{-9}, 1 + 10^{-9}]$ 之间。
\end{itemize}

% 其中至少有 $\left\lfloor \frac{1}{2}n\log_2 n - n \right\rfloor$ 对点对的距离为 $1$。

\subsection{输入格式}

第一行输入一个整数 $T$ ($1 \le T \le 100$)，表示测试数据的数量。

每组测试数据输入一行一个整数 $n$ ($2 \le n \le 5000$)，表示需要选定的点的数量。

保证所有测试数据的 $\sum n \le 5000$。

\subsection{输出格式}

对于每组测试数据，输出 $n$ 行，每行输出两个实数 $x_i, y_i$ ($-10^5 \le x_i, y_i \le 10^5$)，表示第 $i$ 个点的坐标。

你的输出将被视为正确，当且仅当你的输出符合上面的三个条件。

如果有多个合法的答案，则你可以输出其中任意一个。

\subsection{样例 1}

\begin{minipage}{\linewidth}
	\begin{minipage}[t]{0.49\linewidth}
		\begin{lstlisting}[title=\texttt{standard input}]
1
3
\end{lstlisting}
	\end{minipage}
	\begin{minipage}[t]{0.49\linewidth}
		\begin{lstlisting}[title=\texttt{standard output}]
0.0000000000 0.0000000000
1.0000000000 0.0000000000
0.5000000000 0.8660254038
\end{lstlisting}
	\end{minipage}
\end{minipage}


\subsection{样例 1 解释}

对于第一组测试数据，前两个条件不难验证，对于第三个条件，有：

\[d_{12} = 1, d_{13} = d_{23} = 1.000000000013476527219909191607\ldots \approx 1 + 1.3476 \times 10^{-11} \]

因此每个点对之间的距离都在 $1$ 的误差范围内，这是符合第三个条件的。

\subsection{提示}

当 $n = 5000$ 时 $\left\lfloor \frac{1}{2}n\log_2 n - n \right\rfloor = 25719$。

\newpage
\section[E. 占位标题 E]{\myboxv{      \BF{\TT{E}}} 根除}

\textit{运行时间限制}   X second

\textit{内存使用限制} X Gigabyte

\textit{读入文件} \qquad 标准输入

\textit{输出文件} \qquad 标准输出


\begin{quote}
	\textit{占位风味文本}
\end{quote}

\subsection{题目描述}

我们知道对于任意的正整数 $n$ 一定可以写成 $n = ab^2$，其中 $a, b \in \mathbb{N}$。现在定义数论函数 $f(n) = \min\{a \mid n = ab^2, a, b \in \mathbb{N}\}$，即函数值为满足等式的最小正整数 $a$。

现在给定 $q$ 次询问，每次询问给定 $l, r$，请输出 $\sum_{i=l}^r f(i)$ 的答案在模 $998\,244\,353$ 下的值。

\subsection{输入格式}

第一行输入一个整数 $T$ ($1 \le T \le 10^6$)，表示测试数据的数量。

每组测试数据输入一行两个整数 $l$ 和 $r$ ($1 \le l \le r \le 10^{12}$)，表示一次询问的区间。

保证至多有 $10$ 组测试数据的 $r > 10^7$。

\subsection{输出格式}

对于每组测试数据，输出一行一个整数，表示 $\sum_{i=l}^r f(i)$ 在模 $998\,244\,353$ 下的值。

\subsection{样例 1}

\begin{minipage}{\linewidth}
	\begin{minipage}[t]{0.49\linewidth}
		\begin{lstlisting}[title=\texttt{standard input}]
占位输入
\end{lstlisting}
	\end{minipage}
	\begin{minipage}[t]{0.49\linewidth}
		\begin{lstlisting}[title=\texttt{standard output}]
占位输出
\end{lstlisting}
	\end{minipage}
\end{minipage}


\subsection{样例 1 解释}

占位样例解释

\newpage
\section[F. 凸包]{\myboxvi{      \BF{\TT{F}}} 凸包}

\textit{运行时间限制}   2 seconds

\textit{内存使用限制} 1 Gigabyte

\textit{读入文件} \qquad 标准输入

\textit{输出文件} \qquad 标准输出


\begin{quote}
	\textit{我知道既不是凸包也不是凹包的我为什么要哭的。因为我其实是一只豆包。}
\end{quote}

\subsection{题目描述}

给定 $n$ 个一元二次不等式，每个不等式的形式为：

$$
a_i x^2 + b_i x + c_i  >  0
$$

对任意 $1 \le i \le n$ ，保证 $a_i < 0$ 。

请你判断是否存在\BF{整数} $x$，使得这 $n$ 个不等式同时成立。

\subsection{输入格式}

第一行输入一个整数 $T$ ($1 \le T \le 10^4$)，表示测试数据的数量。

每组测试数据第一行输入一个整数 $n$ ($1 \le n \le 2 \times 10^5$)，表示不等式的个数。

随后 $n$ 行，每行输入三个整数 $a_i, b_i, c_i$ ($-10^5 \le a_i \le -1$, $-10^5 \le b_i, c_i \le 10^5$)，表示第 $i$ 个不等式的三个系数。

保证所有测试数据的 $\sum n \le 2 \times 10^5$。

\subsection{输出格式}

对于每组测试数据，若存在整数 $x$ 同时满足所有不等式，输出一行 \TT{YES}，否则输出一行 \TT{NO}。

注意评测时不会区分 \TT{YES} 和 \TT{NO} 的大小写。换言之，当答案是肯定的时候输出 \TT{yes}、\TT{YES}、\TT{Yes}、\TT{YeS} 等都会被认为是正确的。

\subsection{样例 1}

\begin{minipage}{\linewidth}
	\begin{minipage}[t]{0.49\linewidth}
		\begin{lstlisting}[title=\texttt{standard input}]
5
2
-1 0 4
-1 4 -3
3
-1 100000 0
-1 202 10000
-2 400 19998
2
-1 2 0
-1 0 -1
2
-100000 0 100000
-1 0 4
2
-1 2 0
-1 4 -3

\end{lstlisting}
	\end{minipage}
	\begin{minipage}[t]{0.49\linewidth}
		\begin{lstlisting}[title=\texttt{standard output}]
NO
YES
NO
YES
NO
\end{lstlisting}
	\end{minipage}
\end{minipage}


\subsection{样例 1 解释}

对于第 $4$ 组测试数据，当 $x=0$ 时，两个不等式均成立，因此存在整数解。

\newpage
\section[G. 星之彼端]{\myboxvii{      \BF{\TT{G}}} 星之彼端}

\textit{运行时间限制}   8 seconds

\textit{内存使用限制} 1 Gigabyte

\textit{读入文件} \qquad 标准输入

\textit{输出文件} \qquad 标准输出


\begin{quote}
	\textit{今夜天空深远宁静，我们望的是同一颗星星吗？}
\end{quote}

\BF{这是一道通信题，请注意本题不同寻常的数据范围约束。}

\subsection{题目描述}

那是整个ECUST都早已遗忘的久远过去……

在那个无人记得的年代，竹林和千外并肩战斗，共同生活，一起仰望同一片星空。后来，语言被风沙掩埋，道路被海洋截断，连名字都从史书里褪色。唯一没有被带走的，是头顶那片沉默而辽阔的星图。它公开、闪耀、亘古不变，却没有人知道其中哪几颗星曾被他们悄悄连成一句话。

本题共有 $20$ 个测试点。

给定一个无向图 $G = (V, E)$，其中 $|V| = 2000$, $|E| = 43\,600$。图 $G$ 不包含重边和自环，并且保证由以下方式生成：初始时，图 $G$ 包含 $|V|$ 个点且没有任何边。随后，生成器将重复执行以下操作 $|E|$ 次：在当前所有尚未被边连接的点对中，等概率地随机选取一对点，并在它们之间添加一条无向边。

你的程序将被运行两次：

\BF{第一次运行}

在第一次运行时，你将扮演竹林。你需要在图 $G$ 中选择 $n$ 个诱导四元环并输出。这些诱导四元环可以拥有共同的顶点。

\begin{itemize}
	\item \BF{诱导四元环的定义：}选定图 $G$ 中的 $4$ 个点，这 $4$ 个点之间\BF{恰好}存在 $4$ 条边，且构成一个简单的环（即内部不能有对角线相连）。
\end{itemize}

\BF{第二次运行}

在第二次运行时，你将扮演千外。评测机将获取你在第一次运行中输出的所有四元环。你需要回答 $q$ 次询问。对于每一个询问，评测机会挑选一个四元环，任选其中一个点并删除，然后输出剩余三个点的编号。你需要回答这三个点来自于竹林输出的第几个四元环。

\subsection{输入格式}

\BF{第一次运行}

第一行包含一个字符串 \TT{send}。

第二行包含一个整数 $n$（$n = 200\,000$），表示需要输出的诱导四元环数量。

第三行包含两个整数 $|V|$ 和 $|E|$。

接下来 $|E|$ 行，每行包含两个整数 $u, v$，表示顶点 $u$ 和 $v$ 之间有一条无向边 ($1 \le u, v \le |V|$)。

\BF{第二次运行}

第一行包含一个字符串 \TT{receive}。

第二行包含两个整数 $|V|$ 和 $|E|$。

接下来 $|E|$ 行，每行包含两个整数 $u, v$，表示顶点 $u$ 和 $v$ 之间有一条无向边 ($1 \le u, v \le |V|$)。

第 $|E| + 3$ 行包含一个整数 $q$ ($q = 4n$)，表示接收到的顶点集合数量。

接下来的 $q$ 行，每行包含 $3$ 个整数，代表评测机从你第一次运行输出的某个四元环中删除一个顶点后得到的顶点集合。

\subsection{输出格式}

\BF{第一次运行}

输出 $n$ 行。第 $i$ 行输出 $4$ 个整数 $v_1, v_2, v_3, v_4$，代表第 $i$ 个诱导四元环的四个顶点编号。

\BF{第二次运行}

输出 $q$ 行。对于每一组接收到的顶点集合，输出一个整数 $i$ ($1 \le i \le n$)，表示该集合属于第 $i$ 个四元环。

\subsection{样例 1}

\BF{第一次运行}

\begin{minipage}{\linewidth}
	\begin{minipage}[t]{0.49\linewidth}
		\begin{lstlisting}[title=\texttt{standard input}]
send
200000
2000 43600
29 1973
228 1127
...
92 573
\end{lstlisting}
	\end{minipage}
	\begin{minipage}[t]{0.49\linewidth}
		\begin{lstlisting}[title=\texttt{standard output}]
1 174 613 1323
1 955 1683 1755
...
1390 1615 1683 1777
\end{lstlisting}
	\end{minipage}
\end{minipage}

\BF{第二次运行}

\begin{minipage}{\linewidth}
	\begin{minipage}[t]{0.49\linewidth}
		\begin{lstlisting}[title=\texttt{standard input}]
receive
2000 43600
29 1973
228 1127
...
92 573
800000
46 1297 1949
1056 1275 1454
...
387 1289 1533
\end{lstlisting}
	\end{minipage}
	\begin{minipage}[t]{0.49\linewidth}
		\begin{lstlisting}[title=\texttt{standard output}]
17713
191947
...
118231
\end{lstlisting}
	\end{minipage}
\end{minipage}


\subsection{样例 1 解释}

由于完整数据过大，题面省略了样例的部分内容。



\newpage
\section[H. 一般攻击魔法]{\myboxviii{      \BF{\TT{H}}} 一般攻击魔法}

\textit{运行时间限制}   X second

\textit{内存使用限制} X Gigabyte

\textit{读入文件} \qquad 标准输入

\textit{输出文件} \qquad 标准输出


\begin{quote}
	\textit{"Zoltraak."}
\end{quote}

\subsection{题目描述}

在一个 $n \times m$ 的二维网格图中，网格中的每个格子可以处于以下三种状态之一：

\begin{enumerate}
	\item \TT{/}（正斜杠镜子）：双面反光。光线在此发生 $90^\circ$ 偏折，即连通（上方与右方）以及（下方与左方）。
	\item \TT{\char`\\}（反斜杠镜子）：双面反光。光线在此发生 $90^\circ$ 偏折，即连通（上方与左方）以及（下方与右方）。
	\item \TT{.}（空格子）：无障碍物。光线在此沿直线穿透，不发生偏折，即连通（上方与下方）以及（左方与右方）。
\end{enumerate}

初始时，网格中的部分格子已经固定放入了 \TT{/} 或 \TT{\char`\\} 镜子。其余格子目前为空。

对于每一个初始为空的格子，你可以选择\BF{将其保留为空}，或者将其\BF{替换为 \TT{/} 镜子}，或者将其\BF{替换为 \TT{\char`\\} 镜子}。

网格的最外围边界共有 $2(n + m)$ 个"端口"，光线可以从这些边界端口射入或射出网格。

现给定一束光线的\BF{初始射入端口}（及射入方向）以及期望的\BF{目标射出端口}（及射出方向）。

请问是否存在一种合法的网格填充方案，使得光线从指定的射入端口进入后，在网格内经过若干次直线传播与镜面反射，最终\BF{恰好}从指定的目标射出端口射出？

如果存在合法方案，请输出填充后的网格；如果存在多种方案，输出任意一种即可。如果不可能做到，请报告无解。

\subsection{输入格式}

第一行输入一个整数 $T$ ($1 \le T \le 100$)，表示测试数据的数量。

每组测试数据第一行输入两个整数 $n$ 和 $m$ ($1 \le n \times m \le 2 \times 10^5$)，表示网格的行数和列数。

随后 $n$ 行，每行输入一个长度为 $m$ 的字符串，表示网格的初始状态。其中 \TT{/} 表示正斜杠镜子，\TT{\char`\\} 表示反斜杠镜子，\TT{.} 表示空格子。

随后一行输入六个整数 $x_1, y_1, d_1, x_2, y_2, d_2$，其中：

\begin{itemize}
	\item $(x_1, y_1)$ 表示光线射入端口的位置，$d_1 \in \{0, 1, 2, 3\}$ 表示射入方向（$0$ 表示向右，$1$ 表示向下，$2$ 表示向左，$3$ 表示向上）。
	\item $(x_2, y_2)$ 表示光线期望射出端口的位置，$d_2 \in \{0, 1, 2, 3\}$ 表示射出方向（定义同上）。
\end{itemize}

端口坐标的定义为：$(0, j)$ 表示第 $j$ 列的上边界，$(n + 1, j)$ 表示第 $j$ 列的下边界，$(i, 0)$ 表示第 $i$ 行的左边界，$(i, m + 1)$ 表示第 $i$ 行的右边界。保证输入的端口在边界上，且射入方向指向网格内部，射出方向指向网格外部。

保证所有测试数据的 $\sum (n \times m) \le 2 \times 10^5$ 。

\subsection{输出格式}

对于每组测试数据,若存在合法方案，输出 $n$ 行，每行输出一个长度为 $m$ 的字符串，表示填充后的网格。否则，输出一行 \TT{IMPOSSIBLE}。


\subsection{样例 1}

\begin{minipage}{\linewidth}
	\begin{minipage}[t]{0.49\linewidth}
		\begin{lstlisting}[title=\texttt{standard input}]
1
2 2
..
..
1 0 0 0 2 3
\end{lstlisting}
	\end{minipage}
	\begin{minipage}[t]{0.49\linewidth}
		\begin{lstlisting}[title=\texttt{standard output}]
./
\.
\end{lstlisting}
	\end{minipage}
\end{minipage}


\subsection{样例 1 解释}

对于第 $1$ 组测试数据，光线从第 $1$ 行的左边界射入（向右），经过 $(1, 1)$ 遇到 \TT{.} 继续向右，经过 $(1, 2)$ 遇到 \TT{/} 向上偏折，从第 $2$ 列的上边界射出（向上）。

\newpage
\section[I. 天外来物]{\myboxix{      \BF{\TT{I}}} 天外来物}

\textit{运行时间限制}   1 second

\textit{内存使用限制} 1 Gigabyte

\textit{读入文件} \qquad 标准输入

\textit{输出文件} \qquad 标准输出


\begin{quote}
	\textit{在燥热的季节，一颗流星划过天际，留下光亮的拖尾。}
	
	\textit{我循着光奔跑的那个夜晚，便是故事的开端。}
\end{quote}

\subsection{题目描述}

有 $n$ 颗水晶，从左到右编号为 $1$ 到 $n$ 。第 $i$ 颗水晶初始时含有 $a_i$ 单位的能量。你可以按照任意顺序依次采集这 $n$ 颗水晶。

当你采集编号为 $i$ 的水晶时：

\begin{itemize}
	\item 你获得该水晶当前所含的能量值，并累加到总得分中。
	\item 如果编号为 $i - 1$ 和 $i + 1$ 的水晶存在且还未被采集，它（们）的能量值会变成原来的相反数。
\end{itemize}

注意，水晶的编号不会随采集发生变化。

采集完所有 $n$ 颗水晶后，可能的最大总得分是多少？

\subsection{输入格式}

第一行输入一个整数 $T$ ($1 \le T \le 10^4$)，表示测试数据的数量。

每组测试数据第一行输入一个整数 $n$ ($1 \le n \le 2 \times 10^5$)，表示水晶的数量。

第二行输入 $n$ 个整数 $a_1, a_2, \ldots, a_n$ ($-10^9 \le a_i \le 10^9$)，表示每颗水晶初始的能量值。

保证所有测试数据的 $\sum n \le 2 \times 10^5$。

\subsection{输出格式}

对于每组测试数据，输出一行一个整数，表示可能的最大总得分。

\subsection{样例 1}

\begin{minipage}{\linewidth}
	\begin{minipage}[t]{0.49\linewidth}
		\begin{lstlisting}[title=\texttt{standard input}]
2
5
1 2 3 4 5
9
9 -9 8 -2 4 -4 3 -5 3
\end{lstlisting}
	\end{minipage}
	\begin{minipage}[t]{0.49\linewidth}
		\begin{lstlisting}[title=\texttt{standard output}]
15
47
\end{lstlisting}
	\end{minipage}
\end{minipage}

\subsection{样例 1 解释}

对于第 1 组测试数据，按照编号 $3 \to 5\to 1\to 2\to 4$ 的顺序采集可以使得最大得分达到 $15$ 。可以证明这是最大值。

\newpage
\section[J. 时间谜题 2]{\myboxx{      \BF{\TT{J}}} 时间谜题 2}

\textit{运行时间限制}   1 second

\textit{内存使用限制} 256 Megabytes

\textit{读入文件} \qquad 标准输入

\textit{输出文件} \qquad 标准输出


\begin{quote}
	\textit{"对无数时间线建立可持久化字典树即可实现对历史的版本控制"，一位青铜龙朋友总是这么告诉我。我对此深信不疑，以至于忘记了自己倒数的心跳。}
\end{quote}

\subsection{题目描述}

星鱼为了证明他上课从不迟到，坚称自己是在 $7{:}110$ 踏进了 $8{:}00$ 开始上课的教室。这启发我们去定义一种"广义时间表示法"。

在广义时间表示法中，时间被记为 $A{:}B$（$A, B$ 均为非负整数）。

我们定义广义时间之间存在两种\BF{等价关系}：

\begin{enumerate}
	\item 对于任意 $B \ge 60$，$A{:}B$ 与 $(A + 1){:}(B - 60)$ 是表示同一个时刻的等价时间。
	\item 对于任意 $A \ge 24$，$A{:}B$ 与 $(A - 24){:}B$ 是表示同一个时刻的等价时间。
\end{enumerate}

我们称一个时间 $h{:}m$ 是"标准时间"，当且仅当它满足 $0 \le h < 24$ 且 $0 \le m < 60$。

现在给定一个广义时间，请你求出与它等价的唯一"标准时间"。

\subsection{输入格式}

第一行输入一个整数 $T$ ($1 \le T \le 10^5$)，表示测试数据的数量。

每组测试数据输入一行两个整数 $A$ 和 $B$ ($0 \le A, B \le 10^9$)，表示一个广义时间 $A{:}B$。

\subsection{输出格式}

对于每组测试数据，输出一行两个整数 $h$ 和 $m$，表示与给定广义时间等价的标准时间 $h{:}m$。

\subsection{样例 1}

\begin{minipage}{\linewidth}
	\begin{minipage}[t]{0.49\linewidth}
		\begin{lstlisting}[title=\texttt{standard input}]
5
7 110
25 30
0 0
23 59
100 200
\end{lstlisting}
	\end{minipage}
	\begin{minipage}[t]{0.49\linewidth}
		\begin{lstlisting}[title=\texttt{standard output}]
8 50
1 30
0 0
23 59
7 20
\end{lstlisting}
	\end{minipage}
\end{minipage}


\subsection{样例 1 解释}

对于第 $1$ 组测试数据，$7{:}110 \to 8{:}50$（应用规则 1，$110 = 60 + 50$）。

对于第 $2$ 组测试数据，$25{:}30 \to 1{:}30$（应用规则 2，$25 = 24 + 1$）。

对于第 $3$ 组测试数据，$0{:}0$ 本身就是标准时间。


\newpage
\section[K. 永世沙漏]{\myboxxi{      \BF{\TT{K}}} 永世沙漏}

\textit{运行时间限制}   X second

\textit{内存使用限制} X Gigabyte

\textit{读入文件} \qquad 标准输入

\textit{输出文件} \qquad 标准输出


\begin{quote}
	\textit{占位风味文本}
\end{quote}

\subsection{题目描述}

给定一个按照逆时针顺序给出的多边形 $P = P_0P_1 \cdots P_{n-1}$。记下面的一个方向 $d(\alpha) = (\cos \alpha, \sin \alpha)$，$0 \le \alpha < \frac{\pi}{2}$ 为一个\BF{稳定方向}，当且仅当存在一条直线 $l$ 满足：

\begin{itemize}
	\item 其方向向量 $\vec{l}$ 满足 $\vec{l} \cdot d(\alpha) = 0$。
	\item 对于 $P$ 闭集内每一点 $P^*$，过其作垂线 $P^*H \perp l$ 交 $l$ 于 $H$，满足垂线段 $P^*H$ 在 $P$ 闭集内。
\end{itemize}

请求出满足条件的 $\alpha$ 的数量。

\subsection{输入格式}

第一行输入一个整数 $T$ ($1 \le T \le 100$)，表示测试数据的数量。

每组测试数据第一行输入一个整数 $n$ ($3 \le n \le 10^5$)，表示多边形的顶点数量。

随后 $n$ 行，每行输入两个整数 $x_i, y_i$ ($-10^9 \le x_i, y_i \le 10^9$)，表示多边形的第 $i$ 个顶点的坐标，按逆时针顺序给出。

保证多边形是简单多边形（边不相交），且所有测试数据的 $\sum n \le 10^5$。

\subsection{输出格式}

对于每组测试数据，输出一行一个整数，表示满足条件的稳定方向的数量。

\subsection{样例 1}

\begin{minipage}{\linewidth}
	\begin{minipage}[t]{0.49\linewidth}
		\begin{lstlisting}[title=\texttt{standard input}]
占位输入
\end{lstlisting}
	\end{minipage}
	\begin{minipage}[t]{0.49\linewidth}
		\begin{lstlisting}[title=\texttt{standard output}]
占位输出
\end{lstlisting}
	\end{minipage}
\end{minipage}


\subsection{样例 1 解释}

占位样例解释

\newpage
\section[L. 净化行动 3]{\myboxxii{      \BF{\TT{L}}} 净化行动 3}

\textit{运行时间限制}   2 seconds

\textit{内存使用限制} 1 Gigabyte

\textit{读入文件} \qquad 标准输入

\textit{输出文件} \qquad 标准输出


\begin{quote}
	\textit{除恶务尽——何为「恶」，我们自有裁定。}
\end{quote}

\subsection{题目描述}

Maddy 仍在清剿 Baddy 的邪恶军团……（这已经是第三次了！）

游戏在一个无限大的平面直角坐标系（可以看作无限大的棋盘）上进行。Baddy 操控了 $n$ 个\BF{炮}棋子。这些炮分别部署在 $(x_1, y_1), (x_2, y_2), \ldots, (x_n, y_n)$ 上，且固定不会移动。

在象棋规则中，炮的攻击规则是"隔山打牛"。如果两个炮 $A$ 和 $B$ 处于同一行或同一列，且它们之间\BF{有且仅有}一个其他炮 $C$，那么我们称 $A$ 攻击了 $B$（同时 $B$ 也攻击了 $A$），而 $C$ 就是它们攻击的"炮架"。

每一次，Maddy 可以选择一个\BF{当前没有受到任何其他炮攻击}的炮，并将其摧毁。被摧毁的炮会被立即移出棋盘，它将不再能攻击其他炮，也不能再作为其他炮的"炮架"。

Maddy 的目标是摧毁所有的 $n$ 个炮。是否存在一种摧毁顺序，使得 Maddy 可以最终清空棋盘？

\subsection{输入格式}

第一行输入一个整数 $T$ ($1 \le T \le 10^5$)，表示测试数据的数量。

每组测试数据第一行输入一个整数 $n$ ($1 \le n \le 2 \times 10^5$)，表示炮的数量。

随后 $n$ 行，每行输入两个整数 $x_i, y_i$ ($-10^9 \le x_i, y_i \le 10^9$)，表示第 $i$ 个炮的位置。

保证所有测试数据的 $\sum n \le 2 \times 10^5$，且单组测试数据中所有炮的位置两两不同。

\subsection{输出格式}

对于每组测试数据，若存在一种合法的摧毁顺序能够清空棋盘，则输出一行 \TT{YES}，否则输出一行 \TT{NO}。

注意评测时不会区分 \TT{YES} 和 \TT{NO} 的大小写，换言之当答案是肯定的时候输出 \TT{yes}、\TT{YES}、\TT{Yes}、\TT{YeS} 等都会被认为是正确的。

\subsection{样例 1}

\begin{minipage}{\linewidth}
	\begin{minipage}[t]{0.49\linewidth}
		\begin{lstlisting}[title=\texttt{standard input}]
3
3
0 0
0 1
0 2
5
0 0
0 1
0 2
1 1
2 1
4
0 0
1 0
2 0
3 0
\end{lstlisting}
	\end{minipage}
	\begin{minipage}[t]{0.49\linewidth}
		\begin{lstlisting}[title=\texttt{standard output}]
YES
YES
NO
\end{lstlisting}
	\end{minipage}
\end{minipage}


\subsection{样例 1 解释}

对于第 $1$ 组测试数据，三个炮在同一列上。初始时，$(0, 0)$ 和 $(0, 2)$ 互相攻击（以 $(0, 1)$ 为炮架），而 $(0, 1)$ 没有受到攻击。因此可以先摧毁 $(0, 1)$，然后 $(0, 0)$ 和 $(0, 2)$ 都不再受到攻击，可以按任意顺序摧毁。

对于第 $3$ 组测试数据，四个炮在同一行上。初始时，$(0, 0)$ 和 $(2, 0)$ 互相攻击（以 $(1, 0)$ 为炮架），$(1, 0)$ 和 $(3, 0)$ 互相攻击（以 $(2, 0)$ 为炮架）。因此 $(0, 0)$、$(1, 0)$、$(2, 0)$、$(3, 0)$ 都受到攻击，无法摧毁任何一个，因此无解。

\newpage
\section[M. 魔女审判]{\myboxxiii{      \BF{\TT{M}}} 魔女审判}

\textit{运行时间限制}   X second

\textit{内存使用限制} X Gigabyte

\textit{读入文件} \qquad 标准输入

\textit{输出文件} \qquad 标准输出


\begin{quote}
	\textit{占位风味文本}
\end{quote}

\subsection{题目描述}

Celeste 国的航天事业正在蒸蒸日上！

近日，Celeste 国的航天工程师发射了第一个星光站，星光站共有 $n$ 个舱室，编号分别为 $1, 2, \ldots, n$，其中主舱室的编号为 $1$。每个舱室通过若干条双向管道相互连接，保证每个舱室都通过一条唯一的由双向管道组成的路径与主舱室相连（即形成一棵以 $1$ 为根的树），因此双向管道的数量为 $n - 1$。

初始情况下，每个舱室 $i$ 中分别放有一份编号为 $a_i$ 的实验器材，保证所有 $a_i$ 互不相同。

作为星光站主管，Maddy 收到了如下 $q$ 条按顺序发来的指令，她需要根据这些指令依次来进行实验，指令共分为两种：

\begin{enumerate}
	\item \BF{实验器材更新}：指令由两个参数 $i$ 和 $w$ 组成，表示将舱室 $i$ 中的原有实验器材丢弃，并更新为编号为 $w$ 的实验器材。
	\begin{itemize}
		\item 保证 $w$ 不和先前的任意一个 $w$ 相同，亦保证不与任意一个初始的 $a_i$ 相同（即所有出现过的器材编号全局两两不同）。
	\end{itemize}

	\item \BF{实验准备}：指令由两个参数 $l$ 和 $r$ 组成，表示需要寻找一组满足如下条件 $C(l, r)$ 的舱室 $x, y, z$ 并给出回答：
	\begin{itemize}
		\item 所有满足在 $x$ 到 $z$ 路径上，或 $y$ 到 $z$ 路径上的点组成的集合 $V$，恰好是整数区间 $[l, r]$ 内的所有数的一个排列（即集合 $V$ 中的器材编号恰好为 $l, l + 1, \ldots, r$，且每个出现一次）。
		\item $z$ 在 $x$ 到主舱室的路径上，$z$ 在 $y$ 到主舱室的路径上（即 $z$ 是 $x$ 和 $y$ 的某个公共祖先，在根为 $1$ 的树中）。
		\item $x$ 和 $y$，$x$ 和 $z$，以及 $y$ 和 $z$ 都可以相等。
	\end{itemize}
\end{enumerate}

请帮助 Maddy 处理这些指令，并对于每个类型为\BF{实验准备}的指令给出回答。

注意评测时不会区分 \TT{YES} 和 \TT{NO} 的大小写，换言之当答案是肯定的时候输出 \TT{yes}、\TT{YES}、\TT{Yes}、\TT{YeS} 等都会被认为是正确的。

\subsection{输入格式}

第一行输入一个整数 $T$ ($1 \le T \le 10^3$)，表示测试数据的数量。

每组测试数据第一行输入两个整数 $n$ 和 $q$ ($1 \le n, q \le 3 \times 10^5$)，表示星光站的舱室数量与指令数量。

随后 $n - 1$ 行每行输入两个整数 $u$ 和 $v$ ($1 \le u, v \le n$, $u \ne v$)，表示存在一条连接 $u$ 和 $v$ 的双向管道。

随后一行输入 $n$ 个整数 $a_1, a_2, \ldots, a_n$ ($1 \le a_i \le 10^9$)，分别表示舱室初始的实验器材编号。

随后 $q$ 行每行表示一组依次进行的指令，第 $i$ 行首先输入一个字符 $\text{op}_i$（$\text{op}_i \in \{\text{U}, \text{E}\}$），表示指令类型，随后格式如下：

\begin{itemize}
	\item 若 $\text{op}_i = \text{U}$，随后输入两个整数 $i, w$ ($1 \le i \le n$, $1 \le w \le 10^9$)，表示进行一次实验器材更新，将舱室 $i$ 中的原有实验器材丢弃，并更新为编号为 $w$ 的实验器材。
	\item 若 $\text{op}_i = \text{E}$，随后输入两个整数 $l, r$ ($1 \le l \le r \le 10^9$)，表示进行一次实验准备。
\end{itemize}

保证单组测试数据内，所有 $w$ 和所有 $a_i$ 两两不同。

保证所有测试数据的 $\sum n \le 6 \times 10^5$，$\sum q \le 6 \times 10^5$。

\subsection{输出格式}

对于每组测试数据中指令类型为\BF{实验准备}的指令输出一行一个字符串：

\begin{itemize}
	\item 若存在一组如此合法的 $(x, y, z)$，则输出 \TT{YES}；
	\item 反之则输出 \TT{NO}。
\end{itemize}

\subsection{样例 1}

\begin{minipage}{\linewidth}
	\begin{minipage}[t]{0.49\linewidth}
		\begin{lstlisting}[title=\texttt{standard input}]
1
6 7
1 2
1 3
1 4
4 5
4 6
5 3 2 4 6 7
E 4 7
E 3 7
E 3 4
E 3 5
U 5 1
E 1 6
E 1 3
\end{lstlisting}
	\end{minipage}
	\begin{minipage}[t]{0.49\linewidth}
		\begin{lstlisting}[title=\texttt{standard output}]
YES
NO
YES
YES
NO
NO
YES
\end{lstlisting}
	\end{minipage}
\end{minipage}


\subsection{样例 1 解释}

占位样例解释

\newpage
\intentionalblankpage

\end{document}
